% GENERATED FILE. Edit canonical lessons, then run npm run generate:books.
% canonical-source-hash: fnv1a64:1fdf6967aedb5424
% canonical-lessons: ML-C525-narannuka, ML-C525-pottikkarayuka, ML-C525-paribhavam, ML-C525-potticcirikkuka, ML-C525-asayakkulappam

\chapter{Tears and Laughter}
\label{ch:ml-tears-and-laughter}
\hlchaptermodality{525}

\begin{chapteropening}
\textbf{By the end of this chapter:} \emph{I can say to groan, to burst into tears, a hurt feeling, to burst out laughing and confusion.}

Complete the last lesson of chapter 525: I can say to groan, to burst into tears, a hurt feeling, to burst out laughing and confusion.
\end{chapteropening}

\section[\texorpdfstring{\ml{ഞരങ്ങുക} (ñaraṅṅuka)}{ഞരങ്ങുക (ñaraṅṅuka)}]{\ml{ഞരങ്ങുക} (ñaraṅṅuka) — to groan; to moan}

\label{lesson:ML-C525-narannuka}



\begin{quote}
Before the new one: say the Malayalam for to be surprised, then the Malayalam for fretting.
\end{quote}

\subsection*{You'll want to know: \ml{ഞരങ്ങുക}}
\textbf{\ml{ഞരങ്ങുക}} — \emph{ñaraṅṅuka} — \textquotedblleft{}to groan; to moan\textquotedblright{}.

\begin{quote}
\ml{വേദന} \ml{കാരണം} \ml{രോഗി} \ml{രാത്രി} \ml{മുഴുവൻ} \ml{ഞരങ്ങി}.
\end{quote}

\begin{quote}
\emph{vēdana kāraṇaṁ rōgi rātri muḻuvan ñaraṅṅi.}
\end{quote}

\begin{quote}
The patient groaned all night because of the pain.
\end{quote}

Past: \ml{ഞരങ്ങി}, \emph{ñaraṅṅi}.

A low sound of pain; the noun is \ml{ഞരക്കം}.

\begin{cousinweb}
\textbf{\ml{ഞരങ്ങുക}} (\emph{ñaraṅṅuka}) is native Dravidian: DEDR 2904 gives the Malayalam verb as to grumble or groan, with \ml{ഞരക്കം}, a moan (DEDR 2904).
\end{cousinweb}

\begin{grammarlens}[title={What you've built}]
One more word for describing what you lived and felt.
\end{grammarlens}

\subsection*{Your turn}
\begin{itemize}
\raggedright
  \item \emph{Say it:} \emph{ñaraṅṅuka}
  \item \emph{Say it:} \emph{ñaraṅṅuka}, once more
  \item \emph{Say it:} \emph{vēdana kāraṇaṁ rōgi rātri muḻuvan ñaraṅṅi.} — The patient groaned all night because of the pain.
  \item \emph{Recall:} say the Malayalam for to be surprised, then the Malayalam for fretting, then say \emph{ñaraṅṅuka} again
\end{itemize}

\begin{tcolorbox}[breakable,colback=teal!4,colframe=teal!35!black,title={Before you move on}]
What does \ml{ഞരങ്ങുക} mean? (\textquotedblleft{}To groan; to moan\textquotedblright{}.) Say it once more.
\end{tcolorbox}
\section[\texorpdfstring{\ml{പൊട്ടിക്കരയുക} (poṭṭikkarayuka)}{പൊട്ടിക്കരയുക (poṭṭikkarayuka)}]{\ml{പൊട്ടിക്കരയുക} (poṭṭikkarayuka) — to burst into tears}

\label{lesson:ML-C525-pottikkarayuka}



\begin{quote}
Before the new one: say the Malayalam for fretting, then the Malayalam for to groan.
\end{quote}

\subsection*{You'll want to know: \ml{പൊട്ടിക്കരയുക}}
\textbf{\ml{പൊട്ടിക്കരയുക}} — \emph{poṭṭikkarayuka} — \textquotedblleft{}to burst into tears\textquotedblright{}.

\begin{quote}
\ml{സമ്മാനം} \ml{കിട്ടിയപ്പോൾ} \ml{അവൾ} \ml{പൊട്ടിക്കരഞ്ഞു}.
\end{quote}

\begin{quote}
\emph{sammānaṁ kiṭṭiyappōḷ avaḷ poṭṭikkaraññu.}
\end{quote}

\begin{quote}
When she got the prize, she burst into tears.
\end{quote}

Past: \ml{പൊട്ടിക്കരഞ്ഞു}, \emph{poṭṭikkaraññu}.

\ml{പൊട്ടി}, having burst, + \ml{കരയുക}, to cry. The past is \ml{പൊട്ടിക്കരഞ്ഞു}.

\begin{cousinweb}
\textbf{\ml{പൊട്ടിക്കരയുക}} (\emph{poṭṭikkarayuka}) is native Dravidian, a compound of \ml{പൊട്ടുക}, \textquotedblleft{}to burst\textquotedblright{} (DEDR 4490), and \ml{കരയുക}, \textquotedblleft{}to cry\textquotedblright{} (DEDR 1291).
\end{cousinweb}

\begin{grammarlens}[title={What you've built}]
One more word for describing what you lived and felt.
\end{grammarlens}

\subsection*{Your turn}
\begin{itemize}
\raggedright
  \item \emph{Say it:} \emph{poṭṭikkarayuka}
  \item \emph{Say it:} \emph{poṭṭikkarayuka}, once more
  \item \emph{Say it:} \emph{sammānaṁ kiṭṭiyappōḷ avaḷ poṭṭikkaraññu.} — When she got the prize, she burst into tears.
  \item \emph{Recall:} say the Malayalam for fretting, then the Malayalam for to groan, then say \emph{poṭṭikkarayuka} again
\end{itemize}

\begin{tcolorbox}[breakable,colback=teal!4,colframe=teal!35!black,title={Before you move on}]
What does \ml{പൊട്ടിക്കരയുക} mean? (\textquotedblleft{}To burst into tears\textquotedblright{}.) Say it once more.
\end{tcolorbox}
\section[\texorpdfstring{\ml{പരിഭവം} (paribhavaṁ)}{പരിഭവം (paribhavaṁ)}]{\ml{പരിഭവം} (paribhavaṁ) — a hurt feeling; a mild grievance}

\label{lesson:ML-C525-paribhavam}



\begin{quote}
Before the new one: say the Malayalam for to groan, then the Malayalam for to burst into tears.
\end{quote}

\subsection*{You'll want to know: \ml{പരിഭവം}}
\textbf{\ml{പരിഭവം}} — \emph{paribhavaṁ} — \textquotedblleft{}a hurt feeling; a mild grievance\textquotedblright{}.

\begin{quote}
\ml{പിറന്നാളിന്} \ml{വിളിക്കാത്തതിൽ} \ml{സുഹൃത്തിന്} \ml{പരിഭവം} \ml{ഉണ്ട്}.
\end{quote}

\begin{quote}
\emph{piṟannāḷinŭ viḷikkāttatil suhṛttinŭ paribhavaṁ uṇṭŭ.}
\end{quote}

\begin{quote}
My friend is hurt that I didn't call on his birthday.
\end{quote}

A soft, hurt complaint between people close to each other, milder than \ml{ദേഷ്യം}.

\begin{cousinweb}
\textbf{\ml{പരിഭവം}} (\emph{paribhavaṁ}) is borrowed from Sanskrit \emph{paribhava}, \textquotedblleft{}an insult, a humiliation\textquotedblright{} (Monier-Williams). In Malayalam it softened into the hurt one feels at being slighted.
\end{cousinweb}

\begin{grammarlens}[title={What you've built}]
One more word for describing what you lived and felt.
\end{grammarlens}

\subsection*{Your turn}
\begin{itemize}
\raggedright
  \item \emph{Say it:} \emph{paribhavaṁ}
  \item \emph{Say it:} \emph{paribhavaṁ}, once more
  \item \emph{Say it:} \emph{piṟannāḷinŭ viḷikkāttatil suhṛttinŭ paribhavaṁ uṇṭŭ.} — My friend is hurt that I didn't call on his birthday.
  \item \emph{Recall:} say the Malayalam for to groan, then the Malayalam for to burst into tears, then say \emph{paribhavaṁ} again
\end{itemize}

\begin{tcolorbox}[breakable,colback=teal!4,colframe=teal!35!black,title={Before you move on}]
What does \ml{പരിഭവം} mean? (\textquotedblleft{}A hurt feeling; a mild grievance\textquotedblright{}.) Say it once more.
\end{tcolorbox}
\section[\texorpdfstring{\ml{പൊട്ടിച്ചിരിക്കുക} (poṭṭiccirikkuka)}{പൊട്ടിച്ചിരിക്കുക (poṭṭiccirikkuka)}]{\ml{പൊട്ടിച്ചിരിക്കുക} (poṭṭiccirikkuka) — to burst out laughing}

\label{lesson:ML-C525-potticcirikkuka}



\begin{quote}
Before the new one: say the Malayalam for to burst into tears, then the Malayalam for a hurt feeling.
\end{quote}

\subsection*{You'll want to know: \ml{പൊട്ടിച്ചിരിക്കുക}}
\textbf{\ml{പൊട്ടിച്ചിരിക്കുക}} — \emph{poṭṭiccirikkuka} — \textquotedblleft{}to burst out laughing\textquotedblright{}.

\begin{quote}
\ml{തമാശ} \ml{കേട്ട്} \ml{എല്ലാവരും} \ml{പൊട്ടിച്ചിരിച്ചു}.
\end{quote}

\begin{quote}
\emph{tamāśa kēṭṭŭ ellāvaruṁ poṭṭicciriccu.}
\end{quote}

\begin{quote}
Everyone burst out laughing at the joke.
\end{quote}

Past: \ml{പൊട്ടിച്ചിരിച്ചു}, \emph{poṭṭicciriccu}.

\ml{പൊട്ടി} + \ml{ചിരിക്കുക}: laughter that bursts out, like \ml{പൊട്ടിക്കരയുക} for tears.

\begin{cousinweb}
\textbf{\ml{പൊട്ടിച്ചിരിക്കുക}} (\emph{poṭṭiccirikkuka}) is native Dravidian, a compound of \ml{പൊട്ടുക}, \textquotedblleft{}to burst\textquotedblright{} (DEDR 4490), and \ml{ചിരിക്കുക}, \textquotedblleft{}to laugh\textquotedblright{} (DEDR 1562).
\end{cousinweb}

\begin{grammarlens}[title={What you've built}]
One more word for describing what you lived and felt.
\end{grammarlens}

\subsection*{Your turn}
\begin{itemize}
\raggedright
  \item \emph{Say it:} \emph{poṭṭiccirikkuka}
  \item \emph{Say it:} \emph{poṭṭiccirikkuka}, once more
  \item \emph{Say it:} \emph{tamāśa kēṭṭŭ ellāvaruṁ poṭṭicciriccu.} — Everyone burst out laughing at the joke.
  \item \emph{Recall:} say the Malayalam for to burst into tears, then the Malayalam for a hurt feeling, then say \emph{poṭṭiccirikkuka} again
\end{itemize}

\begin{tcolorbox}[breakable,colback=teal!4,colframe=teal!35!black,title={Before you move on}]
What does \ml{പൊട്ടിച്ചിരിക്കുക} mean? (\textquotedblleft{}To burst out laughing\textquotedblright{}.) Say it once more.
\end{tcolorbox}
\section[\texorpdfstring{\ml{ആശയക്കുഴപ്പം} (āśayakkuḻappaṁ)}{ആശയക്കുഴപ്പം (āśayakkuḻappaṁ)}]{\ml{ആശയക്കുഴപ്പം} (āśayakkuḻappaṁ) — confusion; a mix-up}

\label{lesson:ML-C525-asayakkulappam}



\begin{quote}
Before the new one: say the Malayalam for a hurt feeling, then the Malayalam for to burst out laughing.
\end{quote}

\subsection*{You'll want to know: \ml{ആശയക്കുഴപ്പം}}
\textbf{\ml{ആശയക്കുഴപ്പം}} — \emph{āśayakkuḻappaṁ} — \textquotedblleft{}confusion; a mix-up\textquotedblright{}.

\begin{quote}
\ml{പുതിയ} \ml{സമയപ്പട്ടിക} \ml{കാരണം} \ml{കുറച്ച്} \ml{ആശയക്കുഴപ്പം} \ml{ഉണ്ടായി}.
\end{quote}

\begin{quote}
\emph{putiya samayappaṭṭika kāraṇaṁ kuṟaccŭ āśayakkuḻappaṁ uṇṭāyi.}
\end{quote}

\begin{quote}
The new timetable caused some confusion.
\end{quote}

\ml{ആശയം} an idea + \ml{കുഴപ്പം} a tangle; \ml{കുഴപ്പമില്ല}, no problem, holds the same second word.

\begin{cousinweb}
\textbf{\ml{ആശയക്കുഴപ്പം}} (\emph{āśayakkuḻappaṁ}) is a Malayalam compound built from \ml{ആശയം}, from Sanskrit \emph{āśaya} (CDIAL 1453), and \ml{കുഴപ്പം}, \textquotedblleft{}intricacy, trouble\textquotedblright{}, which is native (DEDR 1817).
\end{cousinweb}

\begin{grammarlens}[title={What you've built}]
One more word for describing what you lived and felt.
\end{grammarlens}

\subsection*{Your turn}
\begin{itemize}
\raggedright
  \item \emph{Say it:} \emph{āśayakkuḻappaṁ}
  \item \emph{Say it:} \emph{āśayakkuḻappaṁ}, once more
  \item \emph{Say it:} \emph{putiya samayappaṭṭika kāraṇaṁ kuṟaccŭ āśayakkuḻappaṁ uṇṭāyi.} — The new timetable caused some confusion.
  \item \emph{Recall:} say the Malayalam for a hurt feeling, then the Malayalam for to burst out laughing, then say \emph{āśayakkuḻappaṁ} again
\end{itemize}

\begin{tcolorbox}[breakable,colback=teal!4,colframe=teal!35!black,title={Before you move on}]
What does \ml{ആശയക്കുഴപ്പം} mean? (\textquotedblleft{}Confusion; a mix-up\textquotedblright{}.) Say it once more.
\end{tcolorbox}
